\documentclass[../../main.tex]{subfiles}
\begin{document}
\begin{oframed}
  \subfile{problem}
\end{oframed}

{\bf [解答]}

\begin{figure}[htb]
  \centering
  \begin{tikzpicture}[x=1.9cm, y=1.9cm, >=stealth, every node/.style={font=\small}]
    \coordinate (O) at (0,0);
    \coordinate (One) at (1,0);
    \coordinate (MOne) at (-1,0);
    \coordinate (Beta) at (0,0.55);
    \coordinate (Alpha) at (1.034,1.034);
    \coordinate (AlphaInv) at (-0.484,-0.484);
    % 座標軸
    \draw[->, thick, gray!80!black] (-1.65,0) -- (1.75,0) node[right] {$\mathrm{Re}$};
    \draw[->, thick, gray!80!black] (0,-1.0) -- (0,2.15) node[above] {$\mathrm{Im}$};
    % 円 C (中心 Beta=(0, 0.55), 半径 r = 1.141)
    \draw[line width=1.3pt, blue!85!black] (Beta) circle[radius=1.141];
    % 中心 Beta から各点への半径破線
    \draw[dashed, gray!60] (Beta) -- (One) (Beta) -- (MOne) (Beta) -- (Alpha) (Beta) -- (AlphaInv);
    % 原点を通る直線（弦）
    \draw[dashed, teal!75!black, thick] ($(Alpha)+(45:0.35)$) -- ($(AlphaInv)+(225:0.35)$);
    % 各点
    \node[circle, fill=blue!85!black, inner sep=1.3pt] at (One) {};
    \node[circle, fill=blue!85!black, inner sep=1.3pt] at (MOne) {};
    \node[circle, fill=orange!85!black, inner sep=1.3pt] at (Beta) {};
    \node[circle, fill=gray!80!black, inner sep=1.2pt] at (O) {};
    \node[circle, fill=red!85!black, inner sep=1.4pt] at (Alpha) {};
    \node[circle, fill=red!85!black, inner sep=1.4pt] at (AlphaInv) {};
    % ラベル
    \node[below right=2pt, font=\footnotesize] at (One) {$1$};
    \node[below left=2pt, font=\footnotesize] at (MOne) {$-1$};
    \node[left=3pt, orange!85!black] at (Beta) {$\beta$};
    \node[below right=2pt] at (O) {$O$};
    \node[above right=2pt, red!85!black] at (Alpha) {$\alpha$};
    \node[below left=2pt, red!85!black] at (AlphaInv) {$-\frac{1}{\bar{\alpha}}$};
    \node[above right=3pt, blue!85!black] at (0.9,1.5) {$C$};
  \end{tikzpicture}
  \caption{複素数平面上の円$C$と原点を通る弦の両端点$\alpha,\ -\dfrac{1}{\bar{\alpha}}$}
  \label{fig:complex_circle_inversion}
\end{figure}

$C$が$\pm1$を通るので，$C$の中心は純虚数$\beta$とおける．
$C$の半径を$r\in\mathbb{R}$とすると$C$の方程式は
\begin{align}
  f(z)= |z-\beta|^2=r^2 \label{eq:1}
\end{align}
だから，$z=\pm1, \alpha$を代入して
\begin{align}
  |1-\beta|^2=|-1-\beta|^2=|\alpha-\beta|^2=r^2 \label{eq:2}
\end{align}
である．二つ目の等号から
\begin{align}
  |\beta|^2+(\beta+\bar\beta)+1=|\beta|^2-\bar\alpha\beta-\alpha\bar\beta+|\alpha|^2
\end{align}
だが，$\beta$が準虚数なことより$\beta+\bar\beta=0$だから
\begin{align}
  & |\beta|^2+1=|\beta|^2-\bar\alpha\beta-\alpha\bar\beta+|\alpha|^2 \\
  \therefore
  & |\alpha|^2-\bar\alpha\beta-\alpha\bar\beta-1 = 0 \label{eq:3}
\end{align}
である．

一方で題意より$z=-\dfrac{1}{\bar\alpha}$とすると
\begin{align}
  f(-\dfrac{1}{\bar\alpha})
  &=\left|-\dfrac{1}{\bar\alpha}-\beta\right|^2 \\
  &=|\beta|^2+\dfrac{1}{\alpha}\bar\beta+\dfrac{1}{\bar\alpha}\beta+\dfrac{1}{|\alpha|^2} \\
  &=|\beta|^2+\dfrac{\alpha\bar\beta + \bar\alpha \beta + 1}{|\alpha|^2} \\
  &=|\beta|^2+\dfrac{|\alpha|^2}{|\alpha|^2} \quad (\because\text{\cref{eq:3}}) \\
  &=|\beta|^2+1 \\
  &=r^2 \quad (\because\text{\cref{eq:2}})
\end{align}
である．したがって，$C$は$-\dfrac{1}{\bar\alpha}$も通る．\qed

\newpage

\end{document}
